\documentclass[UTF8,12pt,a4paper,fontset=fandol]{ctexart}

% 支持从项目根目录或当前笔记目录编译。
\makeatletter
\def\input@path{{../}{../../}}
\makeatother
\usepackage{researchnotes}

\title{行列式展开的项数}
\author{}
\date{}

\begin{document}
\maketitle

\section{计数原理与例子}

设 $n$ 为正整数，$A=(a_{ij})$ 为 $n$ 阶实矩阵。
其\keyterm{行列式}（determinant）$\det A$ 完全展开后，
\impo{每个乘积项都从每行、每列各取一个元素，并带有相应的正负号。}

按行依次选取：第一行有 $n$ 列可选，第二行有 $n-1$ 列可选，
依此类推，最后一行只剩一列。每种选法的列指标构成
$1,2,\ldots,n$ 的一个\keyterm{排列}（permutation），
并与展开式中的一项一一对应。因此，由乘法原理，展开项数为
\begin{equation}
  \boxed{n(n-1)\cdots 2\cdot 1=n!}.
  \label{eq:determinant-term-count}
\end{equation}
当 $n=1$ 时，展开式只有一项，也符合 $1!=1$。

\begin{example}
二阶行列式 $\det A=a_{11}a_{22}-a_{12}a_{21}$ 有 $2!=2$ 项；
同理，三阶行列式有 $3!=6$ 项，四阶行列式有 $4!=24$ 项。
\end{example}

这里的 $n!$ 是按元素位置计数的\keyterm{形式展开项数}，
包含可能为零或相互抵消的项，不一定等于代入具体数值、化简后保留的项数。
矩阵本身有 $n^2$ 个元素，这与行列式的展开项数是不同的概念。

\end{document}
